\section{Introduction}

Using LLMs is becoming a practical skill for researchers, builders, and
non-expert users. At the same time, many of the best ways to use LLM agents are
introduced in papers whose methods are difficult to operationalize. A paper may
describe a search policy, memory loop, validation stage, or interface contract,
but a user who wants to reuse the contribution must translate prose into an
agent workflow by hand.

PaperToSkill studies a simple hypothesis: papers can be converted into compact,
editable skills that preserve enough procedural knowledge for agents and users
to reuse the main contribution. We use ``skill'' to mean a natural-language
operational artifact with an entry-point \texttt{SKILL.md}, similar to formats
used by agent harnesses that load task-specific instructions. Unlike a summary,
a skill is not only explanatory. It should specify when to use the method,
which inputs are required, which steps to execute, how to validate progress,
which failures to watch for, and how to adapt the instructions across
harnesses.

This paper takes an artifact-first view. PaperToSkill is not presented as a
fully automatic PDF understanding system. The current implementation converts
curated, source-anchored paper notes into skills and source maps. This narrow
scope is intentional: it lets us evaluate whether the conversion layer
preserves operational knowledge before claiming end-to-end PDF automation.

Our current contributions are: (1) a paper-to-skill schema for operationalizing
research contributions; (2) a deterministic extraction scaffold that emits
\texttt{SKILL.md} plus source maps; (3) deterministic extracted-text-to-note
scaffolds evaluated on Toolformer and AIDE; (4) a four-paper benchmark using AI
Scientist-v2, Reflexion, AIDE, and Toolformer; (5) deterministic/offline
evaluations for structure, context coverage, compactness, source grounding, and
transfer readiness; and (6) a first single-run real-reuse stress test that
records both positive and failed Summary-vs-PaperToSkill outcomes without
promoting them into an aggregate downstream-success claim.
